\begin{abstractCN}
   本文受核桃外壳天然纹理启发，提出一种融合仿生学原理、计算建模、机器学习与智能优化的高性能抗冲击薄壳纹理设计方法。首先，基于图灵斑纹方程生成仿核桃表面特征的随机纹理，通过截面形态表征与起皱机理分析，结合碰撞仿真揭示纹理几何特征对防护性能的影响机制。据此，构建可参数化随机生成的类核桃防护纹理薄壳模型：以规整平面网格为基底，通过纹理特征映射生成三维形貌。对薄壳模型进行撞击模拟，建立面外位移场的响应样本库。针对撞击响应场计算效率瓶颈，提出Mechanics-enhanced图神经网络（GNN）框架。该模型通过学习薄壳几何拓扑与撞击载荷间的隐式物理映射，实现全场变形响应的实时高精度预测。基于上述预测模型，构建进化算法协同的优化设计框架，逆向求解最优纹理空间分布。本研究为仿生抗冲击薄壳结构的高效智能设计提供了新范式。

    \keywordsCN 图神经网络；进化优化；图灵纹理；抗冲击
\end{abstractCN}

\begin{abstractENG}
    This article is inspired by the natural texture of walnut shells and proposes a high-performance impact resistant thin shell texture design method that integrates principles of bionics, computational modeling, machine learning, and intelligent optimization. Firstly, based on the Turing pattern equation, a random texture resembling the surface features of walnuts is generated. Through cross-sectional morphology characterization and wrinkling mechanism analysis, combined with collision simulation, the influence mechanism of texture geometric features on protective performance is revealed. Based on this, a parameterized and randomly generated walnut like protective texture thin shell model is constructed: using a regular planar grid as the basis, a three-dimensional morphology is generated through texture feature mapping. Conduct impact simulation on the thin shell model and establish a response sample library for the out of plane displacement field. A Mechanics enhanced Graph Neural Network (GNN) framework is proposed to address the computational efficiency bottleneck of impact response fields. This model achieves real-time high-precision prediction of the full field deformation response by learning the implicit physical mapping between the geometric topology of the thin shell and the impact load. Based on the above prediction model, a collaborative optimization design framework for evolutionary algorithms is constructed to reverse engineer the optimal texture spatial distribution. This study provides a new paradigm for efficient and intelligent design of biomimetic impact resistant thin shell structures.

    \keywordsENG Graph neural network；sep Evolutionary optimization；sep Turing textures；sep Impact-resistant
\end{abstractENG}
